% TOP_PROBLEM: {"id": 20000923, "title": "Link structure and sharp edge bounds for three-edge tight-path-free triple systems", "category_id": 2, "category": {"id": 2, "name": "combinatorics", "display_name": "Combinatorics", "description": "Counting problems, graph theory, discrete structures.", "slug": "combinatorics", "order_index": 2, "created_at": "2026-07-31T15:26:25.671Z"}, "status": "partially_solved", "problem_number": "AIM-COMBINATORICS-0048", "set_id": 14, "statusQualification": "Defines the main tight-path question; the source discussion of Fano planes and other hypergraphs is contextual, not a replacement target."}
% FIRST_POSTED: 2026-10-01
% ENTRY_KIND: statement_only
% UnsolvedMath ID 20000923; source catalogue code AIM-COMBINATORICS-0048
% !TeX program = lualatex
% Definition and sources only; this file is not an LLM solution attempt.
\documentclass[11pt,a4paper]{article}
\usepackage[margin=25mm]{geometry}
\usepackage{fontspec}
\setmainfont{lmroman10-regular.otf}[BoldFont=lmroman10-bold.otf,
 ItalicFont=lmroman10-italic.otf,BoldItalicFont=lmroman10-bolditalic.otf]
\usepackage{amsmath,amssymb,mathtools}
\usepackage{unicode-math}
\setmathfont{latinmodern-math.otf}
\AtBeginDocument{\let\setminus\smallsetminus}
\usepackage{xurl,hyperref}
\hypersetup{colorlinks=true,urlcolor=blue}
\setcounter{secnumdepth}{0}
\setlength{\parindent}{0pt}
\setlength{\parskip}{5pt}
\setlength{\emergencystretch}{3em}
\begin{document}
\section*{Link structure and sharp edge bounds for three-edge tight-path-free triple systems}\label{20000923}
{\small UnsolvedMath ID 20000923 · AIM-COMBINATORICS-0048 · Combinatorics · AIM Workshop Problem Lists}
\subsection{Definitions and mathematical statement}
A simple $3$-uniform hypergraph has edges that are three-element sets.
Let $P_3$ have vertices $1,2,3,4,5$ and edges
$\{1,2,3\},\{2,3,4\},\{3,4,5\}$. Thus its length is three edges,
and it is a tight path: consecutive edges overlap in two vertices.
Let $K_t^{(3)}$ denote the hypergraph containing all triples on $t$
vertices. Define $R(t)=r_3(P_3,K_t^{(3)})$ to be the least $N$ such
that every red--blue coloring of the triples of $[N]$ contains either
a red copy of $P_3$ or a blue copy of $K_t^{(3)}$.
Copies use distinct vertices and need not be induced.

Determine the asymptotic growth of $R(t)$ as $t\to\infty$.
The workshop specifically asks whether
\[
                         R(t)=o(t^2),
\]
and whether the more precise order $\Theta(t^2/\log t)$ is correct.
Here the first statement means that $R(t)/t^2$ tends to zero; the
second requires positive constant upper and lower multiples for all
sufficiently large $t$.
Equivalently, one seeks the largest number of vertices of a
$P_3$-free triple system whose largest independent set has fewer than
$t$ vertices; an independent set contains no edge. This equivalent
formulation counts the red edges as the hypergraph. The target concerns
exactly three edges, whose intersection contains the middle vertex;
results about longer tight paths do not by themselves settle it.
\subsection{Short English statement}
Determine the Ramsey growth of a three-edge tight triple-system path against a large complete triple system.
\subsection{Sources}
\begin{itemize}
\item[S1] ULAM.AI and the UnsolvedMath Contributors, supplied problems.json, record 20000923 (AIM-COMBINATORICS-0048), AIM Workshop Problem Lists. Catalogue identity and source transcription; CC BY 4.0.
\url{https://huggingface.co/datasets/ulamai/UnsolvedMath}.
\item[S2] American Institute of Mathematics, Graph Ramsey theory, item 1. Original workshop question underlying AIM-COMBINATORICS-0048.
\url{https://www.overleaf.com/read/mnvcscjjysvg}.
\end{itemize}
\end{document}
